% !PRINTED 615
\weberpage{615}

%<<CONT
wird. Es ist dann $c_{1,1}$ die kleinste positive ganze rationale Zahl, die durch $\varphi_1 \varphi_2$ theilbar ist, und diese muss jedenfalls ein Theiler von $a_{1,1} b_{1,1}$ sein. Hat man
\[
a_{1,1} b_{1,1} = k\, c_{1,1},
\]
wo $k$ eine positive ganze rationale Zahl ist, so ist, weil $N_a(\varphi_1)\, N_a(\varphi_2) = N_a(\varphi_3)$ sein muss,
\[
k\, a_{2,2} b_{2,2} = c_{2,2}.
\]
Die letzte Zahl $c_{1,2}$ ist dann aus der Congruenz
\weq{4}{c_{1,2} + c_{2,2}\,\Theta \equiv 0 \ (\mathrm{mod}\ \varphi_1 \varphi_2)}
zu bestimmen.

Wenn es nur darauf ankommt, einen Repr\"asentanten der zusammengesetzten Classe zu finden, so verf\"ahrt man am einfachsten so.

Ist $\varphi_1$ beliebig angenommen, so kann man (nach \S.\,152, 6.) $\varphi_2$ in seiner Classe so w\"ahlen, dass es relativ prim zu $a_{1,1}$ wird. Dann ist auch $\varphi_1$ relativ prim zu $\varphi_2$ und $b_{1,1}$ relativ prim zu $a_{1,1}$.

Denn haben $a_{1,1}$ und $b_{1,1}$ einen gr\"ossten gemeinschaftlichen Theiler $d$, so ist dieser auch relativ prim zu $\varphi_2$, und also ist $b_{1,1} : d$ durch $\varphi_2$ theilbar; dies aber widerspricht, wenn $d > 1$ ist, der Definition von $b_{1,1}$. Ebenso folgt auch umgekehrt, dass, wenn $a_{1,1}$ und $b_{1,1}$ relativ prim sind, $\varphi_2$ zu $a_{1,1}$ relativ prim ist. Denn wenn keiner der rationalen Primfactoren von $b_{1,1}$ in $a_{1,1}$ aufgeht, so kann auch keiner der in $\varphi_2$ aufgehenden Primfactoren in $a_{1,1}$ aufgehen.

In diesem Falle ist nun:
\weq{5}{a_{1,1} b_{1,1} = c_{1,1}, \quad a_{2,2} b_{2,2} = c_{2,2}.}
Die Congruenz (4) wird dann
\weq{6}{c_{1,2} + a_{2,2} b_{2,2}\,\Theta \equiv 0 \ (\mathrm{mod}\ \varphi_1 \varphi_2),}
und da
\[
a_{1,2} + a_{2,2}\,\Theta \equiv 0 \ (\mathrm{mod}\ \varphi_1), \quad b_{1,2} + b_{2,2}\,\Theta \equiv 0 \ (\mathrm{mod}\ \varphi_2)
\]
ist, so zerf\"allt (6) in die beiden Congruenzen:
\[
\begin{aligned}
c_{1,2} &\equiv b_{2,2} a_{1,2} \ (\mathrm{mod}\ \varphi_1) \\
&\equiv a_{2,2} b_{1,2} \ (\mathrm{mod}\ \varphi_2).
\end{aligned}
\]
Diese Congruenzen sind aber nach \S.\,138, 13. gleichwerthig mit den gew\"ohnlichen Zahlencongruenzen
\weq{7}{
\begin{aligned}
c_{1,2} &\equiv b_{2,2} a_{1,2} \ (\mathrm{mod}\ a_{1,1}) \\
&\equiv a_{2,2} b_{1,2} \ (\mathrm{mod}\ b_{1,1}),
\end{aligned}}
wodurch $c_{1,2}$ nach dem Modul $c_{1,1}$ bestimmt ist.
